\section{Parallax correction method}
\label{sec:method}
\label{sec:3}

\subsection{Viewing geometry}\label{sec:3.1}

All positions are geodetic coordinates on the WGS84 ellipsoid \citep{nima2000}. A surface point at latitude \(\varphi\), longitude \(\lambda\) and ellipsoidal height \(h\) has Earth-centred, Earth-fixed (ECEF) coordinates

\[
\mathbf{p} = \begin{pmatrix} (N + h)\cos\varphi\cos\lambda \\ (N + h)\cos\varphi\sin\lambda \\ \left(N(1 - e^2) + h\right)\sin\varphi \end{pmatrix},
\qquad N = \frac{a}{\sqrt{1 - e^2\sin^2\varphi}},
\]

with semi-major axis \(a\) and eccentricity \(e\). The inverse transform uses three fixed-point iterations on latitude. Both directions agree with pyproj \citep{pyproj} to better than 1 mm.

The satellite is at \(\mathbf{s} = R_{\mathrm{geo}}(\cos\lambda_s, \sin\lambda_s, 0)\), with \(R_{\mathrm{geo}} = 42\,164\) km and \(\lambda_s = 74.0^\circ\). The unit line-of-sight vector from the ground to the satellite is \(\hat{\mathbf{d}} = (\mathbf{s} - \mathbf{p})/\lVert\mathbf{s} - \mathbf{p}\rVert\). The viewing zenith angle (VZA) is the angle between \(\hat{\mathbf{d}}\) and the local vertical; the viewing azimuth angle (VAA) comes from its east and north components. These agree with pyorbital \citep{raspaud2018} within 0.01°.

A cloud at height \(h\) appears displaced away from the satellite by approximately

\[
\Delta \approx h \tan(\mathrm{VZA}).
\]

Uncorrected, this displacement misplaces satellite rainfall estimates by several kilometres, especially for deep convection at high viewing angles \citep{vicente2002}. Over India the displacement is mostly northward. For a 12 km cloud top it is about 2.6 km at Kochi (VZA 12°) and 8--9 km over Northeast India and the Himalaya (VZA 34--36°).

\subsection{Cloud-top height}\label{sec:3.2}

Cloud-top height is estimated per pixel from TIR1 brightness temperature. For an opaque cloud, TIR1 BT is close to the cloud-top temperature, so the cloud-top height is the height at which the reference profile has that temperature.

\begin{enumerate}
\def\labelenumi{\arabic{enumi}.}
\tightlist
\item
  \textbf{Cloud mask.} A pixel is cloudy if its BT is more than 8 K below the ISRO clear-sky BT, sampled to the pixel by nearest neighbour. Where no clear-sky BT is available, a pixel is cloudy if BT \textless{} 270 K. Clear pixels get no height and are not moved.
\item
  \textbf{Height lookup.} BT is interpolated linearly in the AFGL tropical profile between the surface and the cold-point tropopause (17 km, 194.8 K). BT colder than the tropopause minimum, as in overshooting tops, is assigned the tropopause height.
\end{enumerate}

As an alternative, ERA5 profiles are interpolated bilinearly to each pixel, with geopotential converted to geometric height. The lookup then runs downwards from each column's cold-point tropopause (searched up to 70 hPa), so low-level inversions cannot capture high clouds. On the 14 validation cases, ERA5 changes cold-core heights by −0.15 km (median) and the INSAT--GPM offsets by under 0.5 km relative to the standard atmosphere.

Profile heights are above mean sea level. We treat them as ellipsoidal heights. The difference (the geoid undulation, −50 to −100 m over India) changes the shift by less than 1\%.

The approach has known limits: semi-transparent cirrus appears warmer than its true top, so its height is underestimated; and the standard atmosphere differs from the real one by a few hundred metres.

\subsection{Ray-tracing solver}\label{sec:3.3}

The navigated (apparent) position of a pixel is where its line of sight meets the ellipsoid. A cloud at height \(h\) lies on the same line of sight, nearer the satellite. The corrected position is therefore the point

\[
\mathbf{q}(t) = \mathbf{p}_0 + t\,\hat{\mathbf{d}}
\]

on the line of sight whose geodetic height equals \(h\). Here \(\mathbf{p}_0\) is the apparent position at the navigation reference height \(h_0\) (\(h_0 = 0\) for L1C data).

We follow \citet{bielinski2020} in two steps:

\begin{enumerate}
\def\labelenumi{\arabic{enumi}.}
\tightlist
\item
  \textbf{Enlarged ellipsoid.} Intersect the line with the ellipsoid of semi-axes \(a + h\) and \(b + h\): \[
  \frac{q_x^2 + q_y^2}{(a + h)^2} + \frac{q_z^2}{(b + h)^2} = 1.
  \] This is a quadratic in \(t\), \(A t^2 + B t + C = 0\). We take the root closest to the apparent point, using the numerically stable form \(t = C/q\) or \(q/A\) with \(q = -\tfrac{1}{2}\left(B + \operatorname{sign}(B)\sqrt{B^2 - 4AC}\right)\). The enlarged ellipsoid is not a surface of constant geodetic height, but the error is only about 1 cm for a 12 km cloud over India.
\item
  \textbf{Newton refinement.} One step on the true geodetic height, \[
  t \leftarrow t + \frac{h - h_{\mathrm{geod}}(\mathbf{q}(t))}{\hat{\mathbf{d}}\cdot\hat{\mathbf{u}}(\mathbf{q}(t))},
  \] where \(\hat{\mathbf{u}}\) is the local vertical at \(\mathbf{q}\). This brings the error below 1 mm.
\end{enumerate}

Points the satellite cannot see (line of sight below the local horizon) and pixels without a height are returned as missing. The corrected position is \(\mathbf{q}(t)\) converted back to geodetic coordinates.

Unlike the spherical-Earth correction in satpy \citep{raspaud2018}, this solver works on the ellipsoid and allows a non-zero reference height. The two agree within 25 m.

\subsection{Terrain}\label{sec:3.4}

Ellipsoid-navigated data misplace not only clouds but also elevated ground: over the Himalaya or the Western Ghats a 2 km ridge is displaced by about 1--1.5 km. With a digital elevation model, the solver:

\begin{itemize}
\tightlist
\item
  \textbf{cloudy pixels:} uses \(\max(h_{\mathrm{cloud}}, h_{\mathrm{ground}})\), so no cloud is placed below the ground;
\item
  \textbf{clear pixels:} moves them to where the line of sight meets the terrain (terrain parallax).
\end{itemize}

The ground height must be taken at the corrected position, which is not known in advance. We solve this by fixed-point iteration (three steps): look up the ground at the current position, correct, and repeat. The iteration converges where the terrain is less steep than the line of sight, which holds for VZA below about 70°.

With SRTM over India, the ground shift reaches 3--4 km over the Himalaya and the Tibetan Plateau, 1--2 km over the Northeast hills and under 0.5 km over the Western Ghats and Deccan. Cloud tops are far above the ground, so terrain changes the position of storm cores negligibly; in the 14 validation cases the offsets with and without terrain are identical to 0.1 km. In clear sky, ground pixels in Himalayan scenes move by 1.3--2.5 km (median), up to about 6 km.

\subsection{Resampling onto a regular grid}\label{sec:3.5}

After correction, pixels no longer lie on a regular grid: several may land in one grid cell and some cells receive none. We resample with three rules:

\begin{enumerate}
\def\labelenumi{\arabic{enumi}.}
\tightlist
\item
  \textbf{Collisions:} when several pixels fall in one cell, the pixel with the greatest height wins. The satellite sees the highest cloud; lower ones beneath it are hidden.
\item
  \textbf{Small gaps:} cells that receive no pixel because neighbouring pixels spread apart are filled by inverse-distance weighting (power 2) from observed cells within one cell. A filled cell stays clear (no cloud-top height) if most of its neighbour weight comes from clear cells.
\item
  \textbf{Large gaps:} cells behind tall clouds, which the satellite never saw, are left empty.
\end{enumerate}

Each output cell carries a flag: 0 observed, 1 interpolated, 2 empty. Later analysis can exclude flags 1 and 2. The default grid is 0.04° (0--40°N, 60--100°E); validation uses 0.02°.

The uncorrected dataset is produced by the same resampling with the apparent positions. Both datasets therefore share exactly the same grid, which is required for the downstream model comparison.

\subsection{Navigation correction}\label{sec:3.6}

Parallax correction assumes the navigated positions are correct. We checked this by matching INSAT-3DR images to coastlines, and found errors of up to 20 km that vary with season.

\textbf{Measurement.} In each of 13 coastal windows around India, Sri Lanka, Bangladesh and Myanmar, we correlate the image with a land-fraction map. The map is made from Natural Earth polygons, blurred to the 4 km pixel footprint.

\begin{itemize}
\tightlist
\item
  \textbf{Day} (solar zenith angle \textless{} 75°): VIS albedo, using clear pixels (VIS \textless{} 30\% and TIR1 \textgreater{} 285 K).
\item
  \textbf{Night:} TIR1 with the sign reversed, since land is colder than sea, using pixels with TIR1 \textgreater{} 285 K.
\end{itemize}

A window is used only if it has at least 300 clear pixels, including at least 50 land and 50 sea pixels. Trial shifts cover ±30 km: first in 2 km steps, then in 0.5 km steps around the best value. The shift with the highest Pearson correlation \(r\) is the window's offset. Two windows whose coasts run roughly north--south or are blurred by estuaries (Chittagong, Sundarbans) measure only the east--west offset.

\textbf{Acceptance.} A window's measurement is kept only if:

\begin{itemize}
\tightlist
\item
  \(r \ge 0.6\), the best shift is not at the search edge, and both components are within 25 km;
\item
  if the window was matched in two or more scans of the same day (for example 04:15 and 06:15 UTC), they agree within 5 km in both components, and their median is used;
\item
  if it was matched in only one scan, \(r \ge 0.75\).
\end{itemize}

These rules reject matches distorted by cloud, which are common in the monsoon.

\textbf{Correction field.} Accepted window offsets are spread over the scene by inverse-distance weighting (power 2, great-circle distance, minimum distance 1 km). The north component uses only windows that measured it. The resulting east and north offsets are added to each pixel's position before cloud-top height and parallax correction.

\textbf{Seasonal behaviour.} Measurements from March 2022 to December 2024 show two states. From late September to mid-March, positions need a 7--20 km westward correction. From mid-April to September the error is small. The switch near the equinoxes is consistent with the satellite's yaw flip. During the monsoon the error is small and cannot be measured reliably for a single day: night-time matching fails because land and sea have nearly the same TIR1 temperature, and daytime matches from scans 1--2 h apart scatter by about 9--10 km. We therefore apply one monsoon correction, the median of daytime matches near the validation storms (36 matches with r ≥ 0.7 within ±25 km): −0.5 km east and +2.5 km north, uncertain by about ±2 km. (The median of all monsoon measurements across India, +1.8 km east and +4.8 km north, is used as a sensitivity test.)

\subsection{Implementation}\label{sec:3.7}

The method is implemented in the open-source Python package \texttt{pygeoparallax} (v0.2.0, MIT licence; https://github.com/bapallinaveen60/PyGeoParallax-IN). The solver has a vectorised NumPy reference and a parallel Numba kernel \citep{lam2015} with identical results. The Numba kernel corrects a full disk in 0.5 s on a laptop (the NumPy version takes about 9 s); a 1-million-pixel India scene, including resampling, takes about 2 s. Unit tests (57) run on Linux and Windows under continuous integration. A command-line interface handles single scenes, parallel batches and navigation measurement.

